% !TeX program = xelatex
% !TeX bibfile = references.bib
% 在本目录运行 latexmk；全文为可编辑草稿，不含真实研究结果。
\documentclass[UTF8,a4paper,zihao=-4,fontset=fandol]{ctexart}
\usepackage{modeling}

% 题目、题数、摘要及附录安排应按实际任务和适用竞赛要求调整。
\newcommand{\papertitle}{\slot{研究对象与核心问题}}

\begin{document}
\begin{center}
  {\songti\bfseries\fontsize{16pt}{20pt}\selectfont\papertitle\par}
  \vspace{1em}
  {\heiti\bfseries\zihao{4} 摘要}
\end{center}

\slot{用一至两句界定研究对象、任务与关键约束，说明需要确定的输出。}

\textbf{针对问题一，}\slot{从决定模型选择的数据特征或机理出发，交代数学表示与求解方法，给出本问关键答案及单位，再填写实际完成的验证、指标和适用条件。}

% 按实际题数复制前一段并改题号；不存在的分问不得保留。
\slot{后续问题应交代相对于前一问新增的因素、继承的输入以及新的答案。摘要数字必须与正文、结果台账一致。}

\slot{用最有证据支持的结论结束，明确必要边界；没有执行的检验不能写成已完成。}

\noindent\textbf{关键词：}\slot{研究对象}；\slot{核心模型}；\slot{求解方法}

\clearpage

\section{问题重述}
\slot{压缩说明题目给定的研究对象、输入数据、时空范围与关键约束。逐问列明需要交付的数值、方案、预测或解释，保留单位和评价口径。}

\section{问题分析}
\slot{逐问说明难点、现实关系向数学问题的转化理由，以及拟采用方法为何适合本题。交代分问间输入输出的依赖，避免重复题目或只堆叠算法名。}

\section{模型假设}
\slot{按“假设内容、依据、简化作用、可能影响”说明真正参与计算的假设。数量由题目决定；不把待检验结论写成假设。}

\section{符号说明}
表\ref{tab:global-symbols}用于集中定义全文共享的变量、单位和下标；局部符号在首次使用处解释。
\begin{table}[htbp]
  \centering
  \caption{主要符号及单位}
  \label{tab:global-symbols}
  \small
  \begin{tabularx}{\linewidth}{lYY}
    \toprule
    符号 & 含义 & 单位或取值域 \\
    \midrule
    \slot{符号} & \slot{实际变量含义} & \slot{单位或集合} \\
    \bottomrule
  \end{tabularx}
\end{table}

% 复制本文件用于后续分问时，修改题号、所有 q1 标签和对应引用。
\section{问题一的模型建立与求解}
\label{sec:q1-model}

\subsection{数据与建模依据}
\slot{说明输入文件、字段、单位、处理规则和处理前后样本量。由数据特征或机理关系解释模型选择；参数说明来源，缺失信息明确标记。}

\subsection{模型建立}
\slot{先说明现实关系，再给出数学表达，随后解释符号、适用域和用途。写清变量、目标、约束、参数来源及边界条件，使公式与代码口径一致。}
% 公式示例见 examples/（默认不加载）；此处只填写本题实际模型。
% 使用 equation/align 环境，并设置 \label{eq:q1-model} 等唯一标签。

\subsection{求解方法与实现}
\slot{说明算法、关键参数、初始条件、停止准则、随机种子、运行环境和脚本入口。区分可行解、启发式最好解与有最优性证明的解，准确报告求解状态。}

\subsection{结果与分析}
\slot{先直接回答本问，再用图表支持答案并解释原因；比较必须具有共同条件和指标。表中数字、单位和口径对应真实结果台账。}
表\ref{tab:q1-results}预留本问核心答案，尚未计算的项目保持待填。
\begin{table}[htbp]
  \centering
  \caption{问题一的核心结果（待填）}
  \label{tab:q1-results}
  \small
  \begin{tabularx}{\linewidth}{lYY}
    \toprule
    输出 & 结果及单位 & 条件与解释 \\
    \midrule
    \slot{输出项} & \slot{真实结果} & \slot{条件与题意} \\
    \bottomrule
  \end{tabularx}
\end{table}

\subsection{模型验证}
\slot{根据本问主要风险选择可行性、量纲守恒、样本外误差、基线比较、扰动或随机误差检查。写明实际执行的方法、范围、结果与证据，不机械套用全部检验。}

% 按真实题数复制 model-q1.tex 为 model-q2.tex 等，再加入 \input。
% 复制时将 q1 标签前缀及相应引用全部改为 q2，避免重复标签。
\section{模型评价与改进}
\textbf{有证据的优点。}\slot{说明具体设计带来的收益，并指出支持它的计算、比较或检验。}

\textbf{主要局限。}\slot{指出未覆盖的条件、潜在偏差，以及它们通过何种路径影响哪些输出。}

\textbf{改进路径。}\slot{对应局限提出需要新增的数据、模型或检验，区分已完成工作与未来建议。}

\section{结论}
\slot{逐问回收答案和适用条件，不引入正文没有的新数字或发现。没有足够证据的判断明确保留。}


% 默认启用 BibTeX。正文或附录中实际 \cite 的条目才会进入文献表。
\bibliographystyle{gbt7714-numerical}
\bibliography{references}

\appendix
\section{结果明细与支撑材料}
\slot{按当前题目及提交规则列出结果明细。指定小表和主要结论应在正文中；完整大表按规定放入附录及支撑材料。允许选录时写明选录范围及完整结果文件名，不以缩小到无法阅读来容纳全部数据。}

\section{支撑文件与复现入口}
\slot{以当前实际文件生成列表，说明功能、输入输出、运行顺序、依赖和已验证命令；不得恢复用户已删除的过程文件。}

\section{程序与计算说明}
\slot{按适用规则附完整源代码或核心部分，明确完整代码所在文件。附录必须与最终支撑材料版本一致；代码和运行截图不得替代正文论证。}

\end{document}
